\section{Compact smoothings and unitary gluing}\label{sec:geometry}

We construct a compact family in which the six node parameters and the
unitary gluing data can be varied independently. The metric is induced
from a fixed smooth ambient K\"ahler metric. In particular its behaviour
near a node differs from that of a hyperbolic collar.

\subsection{Six independent smoothing coordinates}
In $\PP^2$ put
\[
 l_0=x,\qquad l_1=y,\qquad l_2=z,\qquad l_3=x+y+z,
 \qquad F_0=\prod_{i=0}^3l_i.
\]
The four lines have six distinct nodes $p_{ij}$ and no triple
intersection. For an unordered pair $e=\{i,j\}$ choose a nonzero
constant $\gamma_e$ and set
\[
 f_e=\gamma_e\left(\prod_{k\notin e}l_k\right)^2.
\]
Choose local frames of $\mathcal O_{\PP^2}(4)$ near the six nodes
and normalize $f_e(p_e)=1$. At every other node $p_f$, the section
$f_e$ vanishes to order at least two. For $a=(a_e)$ near zero in
$\CC^6$, consider
\begin{equation}\label{eq:quartic-family}
 \cC=\left\{([x:y:z],a):
 F_0+\sum_e a_e f_e=0\right\}\subset\PP^2\times B.
\end{equation}

\begin{proposition}\label{prop:compact-family}
After restricting $B$, the family \eqref{eq:quartic-family} is
proper, projective and flat, with smooth total space. There are
holomorphic base coordinates $\tau=(\tau_e)$ centred at zero and
disjoint node neighbourhoods in which the family is exactly
\begin{equation}\label{eq:node}
 u_ev_e=\tau_e.
\end{equation}
Every fibre with all $\tau_e\ne0$ is a smooth connected plane
quartic of genus three.
\end{proposition}
\begin{proof}
At a central node the spatial Hessian of the local defining
function is nondegenerate. Its derivative in the corresponding
$a_e$ direction is one. Elsewhere on the central fibre its spatial
gradient is nonzero. The total hypersurface is therefore smooth
near the entire central fibre; properness allows one restriction
of the base for which it remains smooth.

Write the defining function near $p_e$ as $f(z,a)$ in two affine
coordinates. The implicit function theorem gives a unique nearby
critical point $z_e(a)$. Its critical value $c_e(a)=f(z_e(a),a)$
satisfies
\[
 \frac{\partial c_e}{\partial a_f}(0)=f_f(p_e)=\delta_{ef}.
\]
The derivative of $z_e$ contributes nothing because the spatial
gradient vanishes. The holomorphic Morse lemma with parameters
then gives $f=c_e(a)+u_ev_e$. One may obtain these coordinates by
translating the critical point, writing the remaining function
using its integral Hessian, and taking a local holomorphic
congruence square root of that nondegenerate symmetric matrix.
This construction depends holomorphically on $a$ and leaves the
base fixed. The inverse function theorem now makes
$\tau_e=-c_e(a)$ simultaneous base coordinates. Restrict to a
polydisc in these coordinates.

The node charts have no singular point in a fibre with all
$\tau_e\ne0$. On the compact region outside them the spatial
gradient remains nonzero. Properness and projectivity follow
from \eqref{eq:quartic-family}. For flatness, its relative
Cartier sequence
\[
 0\longrightarrow\mathcal O_{\PP^2\times B}(-4)
 \longrightarrow\mathcal O_{\PP^2\times B}
 \longrightarrow\mathcal O_{\cC}\longrightarrow0
\]
remains exact after restriction to every base fibre: none of
the defining quartics is the zero polynomial. The local
flatness criterion applies. Finally a smooth plane quartic
is connected and has genus three.
\end{proof}

Give $\PP^2$ its Fubini--Study form normalized so that a line has
area one, and use the restriction to $\cC$ of its sum with a
positive Euclidean base form. The fibre metric is the induced
Fubini--Study metric. Write $dA_\tau$ for its area form and
$C_i$ for the four central lines. Their limiting areas are
$A_i=1$. We retain $A_i$ in formulas to exhibit the mass
normalization; the arguments also apply to a different fixed
smooth ambient K\"ahler metric.

\begin{lemma}\label{lem:ambient}
For sufficiently small fixed node charts, the fibre metric in
\eqref{eq:node} is uniformly comparable to the metric induced
from $\CC^2$. The region $|u_e|,|v_e|<a$ has area at most $Ca^2$,
independently of its annular modulus and of the other small
parameters. On each fixed compact component exterior the
complex structures and metrics converge smoothly.

If $\rho>0$, $b=\max_e|\tau_e|\ll\rho^2$, and the components are
truncated at radius $\rho$, the coefficient errors in a fixed
finite number of derivatives are bounded by $Cb\rho^{-N}$ for
some fixed $N$.
\end{lemma}
\begin{proof}
On a fibre the other base coordinates are constant, and
$v_e=\tau_e/u_e$. The induced Euclidean area is
\[
 \left(1+\frac{|\tau_e|^2}{|u_e|^4}\right)dx\,dy,
 \qquad |\tau_e|/a<|u_e|<a .
\]
Its integral is at most $2\pi a^2$. Smooth positivity of the
ambient metric on the fixed total-space charts gives the
uniform comparison. Away from the node centres the projection
is a submersion; choose smooth component trivializations that
preserve the branch coordinates on collars. Their coefficients
converge smoothly. On a truncated collar, differentiating
$v_e=\tau_e/u_e$ a fixed number of times costs only a fixed
power of $\rho^{-1}$. Smooth dependence on the remaining
parameters contributes $O(b)$ on fixed exteriors. Combining
these estimates proves the last assertion.
\end{proof}

\subsection{The bundle and its actual holonomy}
Fix an integer $r\ge2$. The edge set is that of $K_4$, and
$U_{ij}\in\SU(r)$ denotes transport from the $j$ frame to the
$i$ frame; set $U_{ji}=U_{ij}^{-1}$. This convention fixes the
signs of all later incidence matrices.

\begin{proposition}\label{prop:bundle}
Every tuple $U=(U_e)\in\SU(r)^6$ defines a holomorphic bundle
$E_U$ on a neighbourhood of the entire central fibre in $\cC$,
with a flat unitary connection $D_U$, parallel metric $h_0$,
and a parallel trivialization of its determinant. Its component
restrictions are trivial flat bundles. On a smooth fibre the
holonomy factors through a surjection
\[
 \pi_1(C_\tau)\longrightarrow\pi_1(K_4),
\]
and graph-cycle holonomies are the ordered products of the
specified edge maps.
\end{proposition}
\begin{proof}
Choose fixed $0<a<b$ within the disjoint node charts and
restrict $\max|\tau_e|<a^2$. Let the node open sets be
$|u_e|,|v_e|<b$. Choose pairwise disjoint open component
exteriors $V_i$ whose incidence overlaps are the collars
$a<|u_e|<b$ or $a<|v_e|<b$. Such sets follow from the smooth
exterior trivializations used above. These trivializations
need not be holomorphic: each open subset of the complex
total space has its trivial holomorphic rank-$r$ bundle.
The two branch collars at a node are disjoint because
$|u_ev_e|<a^2$. Thus there are no triple overlaps. Every
incidence overlap is connected, being an annulus times
the restricted base polydisc.

The sets cover the central fibre and, after a further
restriction of the base, every entire fibre. Indeed a
sequence of omitted points over parameters tending to zero
would, by properness, have an omitted limit on the central
fibre. Put constant special unitary transition matrices on
the incidence overlaps, choosing their two values at an
edge so that the induced component transport is $U_{ij}$.
There is no nontrivial triple cocycle condition. The local
connections $d$ and identity metrics agree under these
transitions. They give the asserted bundle, connection,
metric and determinant trivialization.

On a smooth fibre the component pieces and node annuli are
connected, as are their incidence overlaps. A graph walk
is realized by paths within the component pieces and
crossing paths in the annuli. Each closed graph walk
therefore has a lift to a fibre loop. The van Kampen
presentation obtained from this cover, after killing all
patch loops, is the free fundamental group of the
incidence graph, which is the barycentric subdivision of
$K_4$. Patch loops have trivial transport in the specified
flat frames. Their quotient consequently gives the stated
surjection and the ordered transport formula. In particular
each vanishing-cycle loop has trivial holonomy.
\end{proof}

By the unitary correspondence \cite{narasimhan-seshadri},
$E_U|_{C_\tau}$ is polystable of degree zero, and it is stable
exactly when the holonomy is irreducible. These are variations
of holomorphic bundle structure and stabilizer, with fixed
degree. No variation of the ordinary slope on a smooth
curve is involved.

\begin{example}\label{ex:noncommuting}
Take $r=2$, identity transports on $01,02,03,23$, and
\[
 U_{12}=\frac{I+i\sigma_x}{\sqrt2},\qquad
 U_{13}=\frac{I+i\sigma_y}{\sqrt2},
\]
where $\sigma_x,\sigma_y,\sigma_z$ are the Pauli matrices.
The two cycle holonomies have no common complex invariant
line, so every smooth-fibre bundle is stable. They cannot
be removed by component frame changes. In the orthonormal
basis $\sigma_a/\sqrt2$ of $\Herm_0(2)$ their adjoint maps
are
\begin{equation}\label{eq:rotations}
 R_x=\begin{pmatrix}1&0&0\\0&0&1\\0&-1&0\end{pmatrix},
 \qquad
 R_y=\begin{pmatrix}0&0&-1\\0&1&0\\1&0&0\end{pmatrix}.
\end{equation}
Replacing the two holonomies by
$\exp(i\alpha\sigma_x/2)$ and
$\exp(i\beta\sigma_y/2)$ gives irreducible representations
for small nonzero $\alpha,\beta$, reducible ones on either
axis, and trivial holonomy at the origin.
\end{example}

The gluing also has a holomorphic complex-parameter
realization: use holomorphic $\operatorname{SL}(r,\CC)$
incidence matrices on complex parameter neighbourhoods.
The analytic statements below concern the unitary real
locus. They do not assert flatness in additional complex
parameter directions.

\subsection{Scales and the finite quadratic form}
For a smooth fibre set
\begin{equation}\label{eq:weights}
 L_e=\log(1/|\tau_e|),\quad L_*=\min_eL_e,\quad
 w_e=\frac{2\pi}{L_e},\quad w_*=\frac{2\pi}{L_*}.
\end{equation}
The limit is $L_*\to\infty$, with no restriction on
$L_e/L_f$ or on the dependence of $U$ on $\tau$.
The parameters remain fixed during every flow.
On
$\cV=\bigoplus_{i=0}^3\Herm_0(r)$ put
\begin{equation}\label{eq:finite-forms}
 M(X)=\sum_i A_i\tr(X_i^2),\qquad
 Q_{\tau,U}(X)=\sum_{e=ij}w_e
       \|X_i-U_{ij}X_jU_{ij}^*\|_{\HS}^2.
\end{equation}
Each unoriented edge is counted once. Let
$B_{\tau,U}$ be the self-adjoint operator for the inner
product $M$ whose quadratic form is $Q_{\tau,U}$, and
let $\nu_0\le\cdots\le\nu_{m-1}$ be its eigenvalues,
where $m=4(r^2-1)$. Zeros are included.

The kernel consists exactly of compatible tuples.
Transport along a spanning tree determines such a tuple
from its value at one vertex; the remaining edges require
that value to commute with every cycle holonomy. Thus
$\ker B_{\tau,U}$ is the trace-free Hermitian commutant
of the representation. Its dimension may change with
$U$, while the dimension $m$ of the finite space is fixed.
